\chapter{Sécurité}
\label{chap:securite}

\section{Principes}
\begin{itemize}
  \item Deny by default sur tous les flux.
  \item Séparation des zones par VLAN et par firewall.
  \item Isolation stricte des tenants OpenStack.
  \item Séparation des cyberlabs du reste de l'infrastructure.
  \item Principe du moindre privilège pour tous les comptes.
  \item Traçabilité : journal d'audit, centralisation des logs.
\end{itemize}

\section{Identité et accès}
\subsection{Comptes}
\begin{tabularx}{\textwidth}{@{}lX@{}}
\toprule
\textbf{Type de compte} & \textbf{Usage} \\
\midrule
Utilisateur étudiant & Projet individuel, membres de groupes \\
Utilisateur enseignant & Projet individuel + droit de lecture sur les groupes \\
Administrateur système & Accès management, création de projets, modification de quotas \\
Compte de service & Utilisé par l'orchestration pour piloter OpenStack \\
Compte d'urgence & Accès hors bande (iDRAC, console) en cas de perte d'accès \\
\bottomrule
\end{tabularx}

\subsection{Authentification}
\begin{itemize}
  \item Keystone pour OpenStack.
  \item Intégration LDAP/AD de l'école si l'annuaire existe (recommandé).
  \item MFA pour les comptes administrateurs (TOTP via FreeOTP ou équivalent).
  \item Séparation des domaines Keystone : \texttt{default} (admin), \texttt{students}, \texttt{teachers}.
\end{itemize}

\section{Segmentation réseau}
Voir chapitre~\ref{chap:reseau}. Résumé :
\begin{itemize}
  \item Management isolé, accessible uniquement depuis ADMIN ou VPN admin.
  \item Étudiants sur VLAN 60 (VPN) et sur les réseaux tenant OpenStack.
  \item Cyberlabs sur VLAN 80, sans accès Internet par défaut.
  \item DMZ sur VLAN 50, accessible depuis Internet sur ports 80/443 et VPN uniquement.
\end{itemize}

\section{Isolation multi-tenant OpenStack}
\begin{itemize}
  \item Chaque projet OpenStack a son propre réseau Neutron (VXLAN).
  \item Security groups stricts par projet.
  \item Pas de réseau partagé entre projets sauf décision explicite de l'admin.
  \item Floating IPs limitées (1 par projet individuel, 2 par groupe).
  \item Pas d'accès des VMs étudiantes aux réseaux de management ou à l'API OpenStack (sauf Horizon en lecture).
\end{itemize}

\section{Cyberlabs}
\begin{itemize}
  \item Réseau Neutron dédié, isolé du reste.
  \item Pas de sortie Internet par défaut (whitelist si nécessaire).
  \item Pas d'accès aux volumes des autres projets.
  \item Cycle de vie court (création et destruction à la demande).
  \item Traçabilité des actions (audit log dédié).
  \item Supervision en place pour détecter un comportement anormal (scan de ports, tentative d'évasion).
\end{itemize}

\begin{warnbox}[Point de vigilance cyberlabs]
Un environnement volontairement vulnérable peut être utilisé comme point de départ d'une attaque vers le reste de l'infrastructure. Les précautions suivantes sont indispensables :
\begin{itemize}
  \item Aucun flux autorisé entre cyberlab et le reste du cloud.
  \item Aucun flux autorisé entre cyberlab et le management.
  \item Filtrage strict au niveau du SRX.
  \item Limitation du débit et de la durée de vie.
  \item Supervision renforcée.
\end{itemize}
\end{warnbox}

\section{Protection des secrets}
\begin{itemize}
  \item Aucun secret en clair dans Git.
  \item Vault (HashiCorp Vault ou Ansible Vault) pour les credentials de service.
  \item Rotation des mots de passe critiques (admin, service accounts).
  \item Séparation des clés SSH par rôle (pas de clé partagée).
  \item Gestion des secrets OpenStack via Barbican si activé.
\end{itemize}

\section{Durcissement}
\begin{itemize}
  \item Serveurs : SSH par clé uniquement, pas de root direct, fail2ban, mises à jour automatiques de sécurité.
  \item Hyperviseurs : accès iDRAC/iLO sur réseau dédié, mots de passe uniques, firmware à jour.
  \item Switches et SRX : SSH uniquement, désactivation des services inutiles, ACL de management.
  \item OpenStack : API derrière reverse proxy, TLS obligatoire, quotas stricts sur l'API.
\end{itemize}

\section{Supervision et détection}
\begin{itemize}
  \item Centralisation des logs (Loki ou syslog).
  \item Alertes sur les événements critiques : échec d'authentification, modification de quota, création de projet, connexion admin.
  \item Surveillance des flux réseau (NetFlow ou sFlow si disponible sur Cisco).
  \item Surveillance des resources (Prometheus) : CPU, RAM, disque, réseau.
  \item Revue périodique des comptes et des rôles.
\end{itemize}

\section{Sauvegardes et PRA}
Voir chapitre~\ref{chap:exploitation} (Exploitation). Points clés :
\begin{itemize}
  \item Sauvegarde quotidienne des bases (PostgreSQL orchestration, MariaDB OpenStack).
  \item Sauvegarde des configurations réseau (SRX, Cisco).
  \item Sauvegarde des images et volumes critiques.
  \item Test de restauration trimestriel.
  \item PRA : sauvegarde off-site (cloud externe ou site distant).
\end{itemize}